% Ficha de estudio 3 del Taller virtual ATA EPE: mecánica analítica I (ligaduras y grados de libertad).
\documentclass[11pt,a4paper]{article}
% Preámbulo común de las fichas de estudio del Taller virtual ATA EPE.
\usepackage[spanish,es-noshorthands,es-tabla]{babel}
\usepackage[utf8]{inputenc}
\usepackage[T1]{fontenc}
\usepackage{lmodern}
\usepackage[margin=2.2cm]{geometry}
\usepackage[table]{xcolor}
\usepackage{amsmath,amssymb,bm}
\usepackage{tabularx,array,booktabs}
\usepackage{enumitem}
\usepackage{fontawesome5}
\usepackage[most]{tcolorbox}
\usepackage[hidelinks]{hyperref}
\usepackage{listings}

\definecolor{morado}{HTML}{2A1B7E}
\definecolor{amarillo}{HTML}{FFCF2E}
\definecolor{rojo}{HTML}{C4142C}
\definecolor{verde}{HTML}{1F9D55}
\definecolor{gris}{HTML}{5D5A73}

\newcolumntype{Y}{>{\raggedright\arraybackslash}X}
\renewcommand{\arraystretch}{1.25}
\setlist{itemsep=1pt,topsep=3pt}
\numberwithin{equation}{section}
\lstset{basicstyle=\ttfamily\small,columns=fullflexible,keepspaces=true,frame=single,rulecolor=\color{gris!40},
  literate={á}{{\'a}}1 {é}{{\'e}}1 {í}{{\'i}}1 {ó}{{\'o}}1 {ú}{{\'u}}1 {ñ}{{\~n}}1 {Ω}{{$\Omega$}}1
    {θ}{{$\theta$}}1 {λ}{{$\lambda$}}1 {Δ}{{$\Delta$}}1 {μ}{{$\mu$}}1 {·}{{$\cdot$}}1 {←}{{$\leftarrow$}}1 {§}{{\S}}1 {→}{{$\to$}}1 {≠}{{$\neq$}}1 {É}{{\'E}}1 {Á}{{\'A}}1 {∨}{{$\vee$}}1 {∧}{{$\wedge$}}1 {σ}{{$\sigma$}}1 {×}{{$\times$}}1 {≈}{{$\approx$}}1}

\newcommand{\rev}[1]{\widetilde{#1}}
\newcommand{\I}{\mathbf{I}}
\newcommand{\vect}[1]{\bm{#1}}
\newcommand{\A}{\mathcal{A}}

\newtcolorbox{meta}{colback=morado,colframe=morado,coltext=white,arc=3mm}
\newtcolorbox{lectura}{colback=morado!5,colframe=morado,fonttitle=\bfseries,title={\faBook~Lecturas}}
\newtcolorbox{ejercicio}[1]{colback=amarillo!10,colframe=amarillo!80!black,fonttitle=\bfseries,title={\faPencil*~Ejercicio #1}}
\newtcolorbox{control}{colback=rojo!5,colframe=rojo,fonttitle=\bfseries,title={\faQuestionCircle~Preguntas de control}}
\newtcolorbox{codigo}[1][De la ecuación al código]{colback=verde!6,colframe=verde,fonttitle=\bfseries,title={\faCode~#1}}
\newtcolorbox{cierre}{colback=verde!6,colframe=verde,fonttitle=\bfseries,title={\faFlagCheckered~Cierre del módulo}}
\newtcolorbox{nota}[1]{colback=amarillo!12,colframe=amarillo!80!black,fonttitle=\bfseries,title={\faLightbulb~#1}}

% \cabecera{número}{título}{subtítulo}
\newcommand{\cabecera}[3]{%
  \title{\textcolor{morado}{\textbf{Ficha #1. #2}}\\[3pt]\large Taller virtual ATA EPE · #3}
  \author{}\date{}}

\cabecera{3}{Ligaduras, grados de libertad y ejes}{Módulo 3: mecánica analítica I}

\begin{document}
\maketitle
\vspace{-2.5em}

\begin{meta}
\textbf{\textcolor{amarillo}{Meta del módulo.}} Describir un montaje como un sistema con ligaduras, distinguir las
holónomas de las que no lo son, y calcular cuántos grados de libertad tiene y cuáles son: qué movimientos permite y
alrededor de qué ejes. Cara pública: el \textbf{contador de grados de libertad} y el dibujo de los ejes posibles.
\end{meta}

\begin{lectura}
\begin{enumerate}[leftmargin=*]
  \item \textbf{Lanczos, \emph{The Variational Principles of Mechanics}, caps.~I y II}: coordenadas generalizadas,
        ligaduras, desplazamientos virtuales. Es la lectura más formativa del curso.
  \item \textbf{Goldstein, cap.~1} (§1.3--1.4) y \textbf{cap.~2} (§2.4): ligaduras, principio de d'Alembert,
        multiplicadores de Lagrange.
  \item \textbf{Lynch y Park, \emph{Modern Robotics}, cap.~2} (§2.1--2.3): espacio de configuración, fórmula de
        Grübler, ligaduras de Pfaff y holonomía; y §7.1 (cadenas cerradas).
  \item Murray, Li y Sastry, cap.~2 §2.5 y cap.~6: velocidades de cadenas como tornillos.
  \item Ball, \emph{Theory of Screws}, cap.~3: tornillos recíprocos (la base de §4 de esta ficha).
\end{enumerate}
\end{lectura}

% =================================================================================================
\section{Configuración y ligaduras}

Un montaje de $N$ cuerpos libres tiene configuración $(M_1,\dots,M_N)\in SE(3)^N$: $6N$ grados de libertad. Las
uniones restringen esa libertad.
\begin{itemize}[leftmargin=*]
  \item \textbf{Holónomas}: relaciones entre posiciones, $\vect C(\vect q, t) = \vect 0$. Reducen la dimensión del
        espacio de configuración. Ejemplos: una bisagra; dos ruedas en contacto con ejes fijos
        ($R_A\varphi_A + R_B\varphi_B = $ cte).
  \item \textbf{No holónomas}: relaciones entre velocidades que no son derivadas de ninguna relación entre
        posiciones, $\mathbf A(\vect q)\,\dot{\vect q} = \vect 0$ (forma de Pfaff), no integrable. Ejemplo: una rueda
        que rueda sin deslizar sobre la mesa y puede girar de rumbo: llega a cualquier posición, aunque en cada instante
        solo puede moverse hacia adelante.
\end{itemize}

\begin{ejercicio}{3.1}
Para la rueda que rueda sobre el plano (coordenadas $x, y$, rumbo $\phi$, giro $\theta$, radio $R$), escribe las dos
ligaduras de rodadura en forma de Pfaff. Muestra que no son integrables (usa el teorema de Frobenius o construye una
maniobra que cambie $x$ sin cambiar $y$, $\phi$ ni $\theta$ al final).
\end{ejercicio}

% =================================================================================================
\section{Desplazamientos virtuales y multiplicadores}

Un desplazamiento virtual $\delta\vect q$ es compatible con las ligaduras en un instante congelado:
$\mathbf J\,\delta\vect q = 0$, con $\mathbf J = \partial\vect C/\partial\vect q$. El principio de d'Alembert dice que
las fuerzas de ligadura no hacen trabajo virtual. De ahí, las ecuaciones de Lagrange con multiplicadores:
\begin{equation}
  \frac{d}{dt}\frac{\partial\mathcal L}{\partial\dot{\vect q}} - \frac{\partial\mathcal L}{\partial\vect q}
  = \vect Q + \mathbf J^{\mathsf T}\vect\lambda,\qquad \vect C(\vect q) = \vect 0 .
\end{equation}
Los multiplicadores $\vect\lambda$ son, literalmente, las fuerzas que hacen las uniones: en el módulo 4 el núcleo
los calcula para mantener las ligaduras, y en el 7 se muestran.

% =================================================================================================
\section{Contar grados de libertad}

\subsection{Grübler--Kutzbach}
Para $n$ eslabones (incluido el fijo) y $j$ uniones con $f_i$ grados de libertad cada una:
\begin{equation}
  \text{GDL} = 6\,(n - 1 - j) + \sum_{i=1}^{j} f_i\quad\text{(espacio)},\qquad
  \text{GDL} = 3\,(n - 1 - j) + \sum_{i=1}^{j} f_i\quad\text{(plano)}.
\end{equation}

\begin{ejercicio}{3.2}
Aplica las dos fórmulas a los prototipos de la cartilla: la balanza, la biela-manivela (cuatro eslabones), el
pantógrafo y el sistema de dos ruedas. Para el cuadrilátero articulado, la fórmula espacial da $-2$ y la plana da
$1$. ¿Cuál es correcta y por qué falla la otra?
\end{ejercicio}

\subsection{El rango del jacobiano}
La fórmula de Grübler supone que las ligaduras son independientes. Cuando la geometría es especial (ejes paralelos,
como en todo el kit), hay ligaduras redundantes y la fórmula se equivoca. Lo correcto es
\begin{equation}
  \text{GDL} = 6N - \operatorname{rango}\mathbf J,
\end{equation}
calculado numéricamente. Los movimientos admisibles son el núcleo de $\mathbf J$; cada vector del núcleo da un
bivector por cuerpo: el eje alrededor del cual se mueve cada pieza en ese modo.

% =================================================================================================
\section{Ligaduras como tornillos (en PGA)}

En PGA, la velocidad relativa que permite una unión es un subespacio de bivectores:

\begin{tabularx}{\linewidth}{@{}l l Y@{}}
\toprule
\textbf{Unión} & \textbf{Velocidades relativas permitidas} & \textbf{Ligaduras}\\
\midrule
Revolución, eje $L$ & $\{\lambda L\}$ & 5\\
Cilíndrica (tornillo sin tuerca) & $\{\lambda L + \mu\,\I L\}$ & 4\\
Tornillo con paso $h$ (tuerca enroscándose) & $\{\lambda(L + h\,\I L)\}$ & 5\\
Fija (tuerca apretada) & $\{0\}$ & 6\\
\bottomrule
\end{tabularx}

\medskip
Las fuerzas que hace una unión son las que \emph{no hacen trabajo} con los movimientos que permite: los forques $W$ con
potencia virtual nula, $W\vee\Omega = 0$, para toda $\Omega$ permitida. En el lenguaje de Ball son los
\textbf{tornillos recíprocos}. Así, cada unión da sus filas de $\mathbf J$ de forma directa: una base del espacio
recíproco.

\subsection{Cadenas cerradas}
En un lazo cerrado (cuadrilátero, biela-manivela), recorrer el lazo y volver al punto de partida exige que la suma de
las velocidades de las uniones, todas en el marco del espacio, sea nula:
\begin{equation}
  \sum_{i=1}^{j}\dot q_i\,S_i = 0,\qquad\Longrightarrow\qquad \text{GDL} = j - \operatorname{rango}[\,S_1\ \cdots\ S_j\,].
\end{equation}

\begin{ejercicio}{3.3}
En el cuadrilátero plano, los cuatro ejes $S_i$ son rectas paralelas a $z$. Muestra que generan un espacio de dimensión
3 (no 4 ni 6) y que, por tanto, $\text{GDL} = 4 - 3 = 1$. Esta es la explicación de la «paradoja» del ejercicio 3.2.
\end{ejercicio}

% =================================================================================================
\section{Discretización: álgebra lineal propia}

\begin{codigo}
\begin{lstlisting}
rango(J, tolerancia):
    valores ← valoresPropios(J^T J)        # Jacobi: matriz simétrica 6N×6N
    σ ← raíz de cada valor
    devolver cuántos σ > tolerancia · σ_max · tamaño
nucleo(J): vectores propios de J^T J con σ ≈ 0
\end{lstlisting}
El umbral relativo es necesario: con redondeo, «cero» nunca es cero. Para montajes grandes, $\mathbf J^{\mathsf T}\mathbf J$
pierde precisión (eleva al cuadrado el número de condición); si hace falta, se cambia a una descomposición QR con
pivoteo.
\end{codigo}

\begin{codigo}[Qué programar]
\begin{tabularx}{\linewidth}{@{}l Y@{}}
\texttt{nucleo/lineal.js} & Matrices densas: producto, transpuesta, Jacobi para matrices simétricas, rango, núcleo.\\
\texttt{nucleo/tornillos.js} & Base de bivectores, potencia virtual $W\vee\Omega$, complemento recíproco de un
  subespacio.\\
\texttt{analisis/gdl.js} & Jacobiano del montaje a partir de sus uniones; grados de libertad; modos (ejes por
  cuerpo); comparación con Grübler para avisar de redundancias.\\
\texttt{app/} & Indicador «este montaje tiene $n$ grados de libertad» y dibujo de los ejes de cada modo.\\
\end{tabularx}
\end{codigo}

\begin{control}
\begin{enumerate}[leftmargin=*]
  \item ¿Qué diferencia una ligadura holónoma de una no holónoma? Da un ejemplo del kit de cada una.
  \item ¿Por qué las fuerzas de ligadura no hacen trabajo virtual?
  \item ¿Qué representan los multiplicadores de Lagrange?
  \item ¿Por qué falla la fórmula de Grübler en los mecanismos del kit?
  \item ¿Qué es un tornillo recíproco y por qué describe las fuerzas de una unión?
  \item ¿Por qué hace falta una tolerancia relativa para calcular el rango?
\end{enumerate}
\end{control}

\begin{cierre}
Ejercicios 3.1--3.3 · prueba P12 (balanza 1, biela-manivela 1, montaje trabado 0) y comparación con Grübler · cara
pública: el contador de grados de libertad con los ejes dibujados.
\end{cierre}

\end{document}
